\documentclass[11pt,a4paper]{article}
\usepackage[margin=2.4cm]{geometry}
\usepackage{amsmath,amssymb}
\usepackage{enumitem}
\usepackage{hyperref}
\usepackage[T1]{fontenc}
\usepackage[utf8]{inputenc}

\setlist[itemize]{topsep=2pt,itemsep=2pt}
\setlist[enumerate]{topsep=2pt,itemsep=2pt}

\title{Project 2 Intuition Guide\\WaveFEM Solver}
\author{Numerical Methods for PDE}
\date{}

\begin{document}
\maketitle

\section*{Purpose}
This document explains what we are going to build before writing the full codebase.
It is written for computer science students with weak PDE background.

\section{What The Project Is About}
We must implement a finite element solver for the 2D wave equation:
\[
\frac{\partial^2 u}{\partial t^2} - \Delta u = f \quad \text{in } \Omega
\]
with boundary and initial data:
\[
u = g \quad \text{on } \partial\Omega,
\qquad
u(0) = u_0,
\qquad
\partial_t u(0) = u_1.
\]
Intuitively, \(u(x,y,t)\) is a displacement. Think of a vibrating membrane, like a drum surface.
At every point \((x,y)\) and time \(t\), \(u\) tells how high or low the membrane is.

The project requires both implementation and discussion of space discretization, time discretization,
stability, numerical dissipation, numerical dispersion, and computational aspects.

\section{Meaning Of The PDE}
The equation is
\[
u_{tt} - \Delta u = f.
\]
The term \(u_{tt}\) is acceleration in time. The term \(-\Delta u\) describes spatial tension:
it compares the value at one point with nearby values. The term \(f\) is an external force.

The boundary condition \(u=g\) fixes the membrane on the boundary. The easiest first case is
\(g=0\), meaning the boundary is fixed at height zero. The initial condition \(u_0\) gives the
initial shape, while \(u_1\) gives the initial velocity.

\section{Weak Formulation}
The computer does not solve the PDE point by point. FEM solves an integrated version.
We multiply the PDE by a test function \(v\) and integrate:
\[
\int_\Omega u_{tt} v\,dx - \int_\Omega \Delta u\,v\,dx
=
\int_\Omega f v\,dx.
\]
After integrating by parts the Laplacian term:
\[
\int_\Omega u_{tt} v\,dx
+
\int_\Omega \nabla u \cdot \nabla v\,dx
=
\int_\Omega f v\,dx.
\]
For homogeneous Dirichlet boundary conditions, the boundary term disappears because the test
functions are zero on the boundary.

Final weak form:
\[
(u_{tt},v)_\Omega + (\nabla u,\nabla v)_\Omega = (f,v)_\Omega
\qquad \forall v \in V.
\]
The first term becomes the mass matrix part, the second becomes the stiffness matrix part,
and the right-hand side becomes the force vector.

\section{From Weak Form To Matrices}
In FEM, we approximate
\[
u_h(x,y,t) = \sum_j U_j(t)\phi_j(x,y).
\]
The functions \(\phi_j\) are basis functions. The unknowns are the coefficients \(U_j(t)\).
After inserting this approximation into the weak form, we get
\[
M\ddot U(t) + K U(t) = F(t),
\]
where
\[
M_{ij} = \int_\Omega \phi_j\phi_i\,dx,
\qquad
K_{ij} = \int_\Omega \nabla\phi_j \cdot \nabla\phi_i\,dx,
\qquad
F_i = \int_\Omega f\phi_i\,dx.
\]
This is the key translation from mathematics to code.

\section{Time Discretization}
The wave equation is second order in time, so we need displacement and velocity, or current
and previous displacement. A simple centered finite difference method is
\[
M\frac{U^{n+1}-2U^n+U^{n-1}}{\Delta t^2} + K U^n = F^n.
\]
Equivalently,
\[
M U^{n+1}
=
\Delta t^2 F^n
-\Delta t^2 K U^n
+2MU^n
-MU^{n-1}.
\]
This method is intuitive and second order in time, but conditionally stable: \(\Delta t\)
must be small enough compared to the mesh size.

Another option is an implicit Newmark method. It is usually more stable for larger time steps,
but each time step is more expensive because it requires solving a linear system involving
both \(M\) and \(K\).

\section{Stability, Dissipation, Dispersion}
Stability means the numerical solution does not blow up for no physical reason.
Numerical dissipation means the method artificially damps the wave. Numerical dispersion means
waves travel at wrong numerical speeds, so the shape may shift or distort.

Simple checks:
\begin{itemize}
\item with \(f=0\) and \(g=0\), energy should stay roughly constant for non-dissipative methods;
\item smaller \(h\) and smaller \(\Delta t\) should improve accuracy;
\item the wave amplitude should not disappear too quickly unless the method is dissipative;
\item the numerical wave should not arrive too early or too late compared to an exact solution.
\end{itemize}

\section{First Validation Test}
Use the unit square \(\Omega=(0,1)^2\) and the exact solution
\[
u(x,y,t) = \sin(\pi x)\sin(\pi y)\cos(\sqrt{2}\pi t).
\]
It satisfies homogeneous Dirichlet conditions because the sine is zero on the square boundary.
For this solution:
\[
f=0,
\qquad
u_0(x,y)=\sin(\pi x)\sin(\pi y),
\qquad
u_1(x,y)=0.
\]
This is a good first test because we know what the solution should look like.

\section{Implementation Plan}
The code should follow the standard deal.II lab style.
Suggested structure:
\begin{itemize}
\item \texttt{WaveEquation.hpp}: class declaration;
\item \texttt{WaveEquation.cpp}: implementation;
\item \texttt{main.cpp}: parameters and run call;
\item \texttt{CMakeLists.txt}: build configuration;
\item \texttt{README.md}: build/run instructions;
\item \texttt{mesh/}: optional mesh scripts;
\item \texttt{output/}: generated results.
\end{itemize}

Main methods:
\begin{itemize}
\item \texttt{setup()}: create mesh, finite element space, DoFs, matrices, vectors;
\item \texttt{assemble\_matrices()}: assemble \(M\) and \(K\);
\item \texttt{interpolate\_initial\_conditions()}: fill initial displacement and velocity;
\item \texttt{solve\_time\_step()}: compute the next solution;
\item \texttt{output\_results()}: write files for ParaView;
\item \texttt{compute\_errors()}: compare with exact solution when available;
\item \texttt{run()}: orchestrate the simulation.
\end{itemize}

\section{Algorithm In Pseudocode}
\begin{verbatim}
read parameters
create mesh
create finite element space
distribute degrees of freedom
assemble mass matrix M
assemble stiffness matrix K
interpolate u0 and u1
compute U^0
compute U^1 or initialize velocity

for each time step:
    assemble or evaluate F(t)
    compute U^{n+1}
    apply Dirichlet boundary conditions
    compute optional error/energy
    write output every few steps

write summary CSV files
\end{verbatim}

\section{What To Look At In ParaView}
The most important visualization is the displacement field \(u\). Useful outputs are:
final displacement at time \(T\), animation over time, coarse/fine mesh comparison,
and large/small \(\Delta t\) comparison.

Check that the wave moves smoothly, the boundary condition is respected, the solution is not
exploding, and the amplitude is plausible.

\section{Experiments To Include}
Minimum useful experiments:
\begin{enumerate}
\item exact solution test on \((0,1)^2\);
\item error vs mesh size \(h\);
\item error vs time step \(\Delta t\);
\item energy over time for the homogeneous case;
\item visual comparison in ParaView.
\end{enumerate}

\section{Group Split}
One possible split for five people:
\begin{itemize}
\item Person 1: weak formulation and report explanation;
\item Person 2: deal.II mesh, DoFs, matrices;
\item Person 3: time discretization and solver;
\item Person 4: validation, errors, energy plots;
\item Person 5: ParaView output, README, final presentation.
\end{itemize}

\section{Oral Explanation}
Short explanation:
\begin{quote}
We solve the 2D wave equation with the finite element method in space. After multiplying
by test functions and integrating by parts, we obtain a weak form. Using Lagrange basis
functions, this becomes a second-order ODE system in time, \(M U'' + K U = F\).
We then choose a time discretization, apply Dirichlet boundary conditions, and advance
the solution step by step. We validate the solver with an exact sinusoidal solution and
study stability, dissipation, dispersion, and computational cost.
\end{quote}

Even shorter:
\begin{quote}
FEM turns the PDE into \(M U'' + K U = F\). Then the time method turns that ODE system
into linear algebra at each time step.
\end{quote}

\section{Risk Checklist}
\begin{itemize}
\item sign of the Laplacian term: the weak form should have \(+\nabla u\cdot\nabla v\);
\item boundary conditions: homogeneous Dirichlet is the easiest first case;
\item initial velocity: the wave equation needs both \(u_0\) and \(u_1\);
\item time step: explicit methods can blow up if \(\Delta t\) is too large;
\item output frequency: do not write too many huge files;
\item exact solution: verify that \(f\), \(u_0\), \(u_1\), and \(g\) match the formula;
\item convergence: decreasing \(h\) and \(\Delta t\) should improve accuracy;
\item energy: for the homogeneous case, energy should be approximately conserved unless
the method is dissipative.
\end{itemize}

\section{Final Mental Model}
The project is a pipeline:
\[
\begin{aligned}
\text{PDE problem}
&\to
\text{weak formulation}
\to
\text{FEM matrices } M,K \\
&\to
\text{time stepping}
\to
\text{output}
\to
\text{validation and report}.
\end{aligned}
\]
If we understand this pipeline, the implementation becomes a structured programming task
instead of a mysterious PDE problem.

\end{document}
